\section{Notation and extremal hypotheses}\label{sec:paper-notation}
All graphs are finite and nonempty. An oriented graph has no loops and no
pair of oppositely directed arcs. Set
\[
 N_D^{++}(v)=\left(\bigcup_{a\in N_D^+(v)}N_D^+(a)\right)
       \setminus\bigl(N_D^+(v)\cup\{v\}\bigr),
 \qquad g_D(v)=d_D^+(v)-d_D^{++}(v).
\]
Reverse exact second neighbourhoods are denoted by $N_D^{--}(v)$.
We call $D$ all-deficient when $g_D(v)>0$ for every vertex. A missing
partner of $v$ is a different vertex nonadjacent to it; its set and size
are $\mathcal M_D(v)$ and $\mu_D(v)$. The number of missing unordered
pairs is $M(D)$, so $\sum_v\mu_D(v)=2M(D)$.

The far set and the reverse far set are
\[
 U_D(v)=V(D)\setminus\bigl(\{v\}\cup N_D^+(v)\cup N_D^{++}(v)\bigr),
 \qquad P_D(v)=\{x\ne v:v\in U_D(x)\}.
\]
For $S\subseteq V(D)$ we use external boundaries
\[
 \Gamma_D(S)=\bigcup_{x\in S}N_D^+(x),\qquad
 \partial_D S=\Gamma_D(S)\setminus S,\qquad
 \partial_D^2S=\Gamma_D(\partial_D S)\setminus(S\cup\partial_D S).
\]
The word external distinguishes these sets from conventions that allow a
set neighbourhood to meet the starting set.

We will use three different extremal choices. An all-deficient graph is
\emph{arc-set-minimal} when deleting every nonempty set of arcs on its fixed
vertex set destroys all-deficiency. The abbreviation MIN means, in addition,
that it is strongly connected. A \emph{minimum-order} counterexample has the
smallest order among all counterexamples, including graphs obtained by
rewiring. The lexicographic normal form in Section~\ref{sec:paper-predecessor-cloning}
adds further global objectives. These conditions are not interchangeable.
If a counterexample exists, a MIN graph exists by taking an arc-minimal
spanning subgraph and then a sink strong component. A minimum-order
counterexample is automatically strongly connected.

\begin{definition}[Protection]\label{def:paper-protection}
Write $a\triangleleft_D b$ if $a\to b$ and
$N_D^-(a)\subseteq N_D^-(b)\setminus\{a\}$.
Then $a$ is a protected far head of $b$. This is a strict partial order.
Put
\[
 \mathcal F_D(b)=U_D(b)\setminus\{a:a\triangleleft_D b\},
 \qquad t_D(b)=|\mathcal F_D(b)|.
\]
\end{definition}
Indeed, an alleged path $b\to x\to a$ would force $x\to b$, a digon.
Transitivity follows because $a\to b\to c$ and predecessor containment
give $a\to c$ and the required further containment. A whole vertex, meaning
$\mu_D(v)=0$, has every far head protected.

A direction $a\to b$ on a missing pair is \emph{convenient} if every
original predecessor of $a$ already reaches $b$ in one or two steps.
A missing pair is good if at least one direction is convenient, and bad
otherwise. All such predicates carry an explicit graph index whenever the
graph has changed. A \emph{global median order} maximizes the total number
of forward arcs among all vertex orders. Its last vertex is its feed.
Local feedback inequalities alone do not justify the sedimentation step
used later.

\begin{remark}[Two kinds of unit]
Gap one means $g_D(v)=1$. Residual one means $\eta_D(v)=1$.
Actual failure-source status refers to a witnessed nonconvenient direction
of an original bad pair. None of these three notions is identified with
another without an explicit lemma.
\end{remark}
